\section{Non-Functional Requirements}

Non-functional requirements describe how well the system must perform its functions, not what those functions are \cite{sommerville2016se,iso29148}. For a security tool they are especially important. A detector that finds every attack but raises thousands of false alarms, or that cannot explain its alerts, will not be used in practice. The requirements below set expectations for the proposed system. Where a module has already been measured, the current value is stated so that later increments can be checked against it.

\subsection{Software Quality Attributes}

\begin{itemize}
   \setlength{\itemsep}{1pt}
   \setlength{\parskip}{0pt}
   \setlength{\parsep}{0pt}
   \item \textbf{Accuracy:} Each model shall be evaluated on held-out data not used for training or tuning, and shall report class-level precision, recall and F1-score in addition to accuracy, because overall accuracy is misleading on imbalanced security data \cite{saito2015pr}. Rare classes with too few test examples for a reliable estimate shall be identified as such.
   \item \textbf{Explainability:} Every alert or prediction shown to an analyst shall include the evidence behind it, such as confidence values or the most influential or most anomalous features.
   \item \textbf{Usability:} The planned dashboard shall let an analyst move from a summary of alerts to the details of a single alert in a few steps, using wording that does not require machine learning knowledge.
   \item \textbf{Maintainability:} Modules shall be independent, with model settings kept in configuration files, so that one module can be retrained or replaced without changing the others.
   \item \textbf{Compatibility:} The system shall run on Windows and Linux with a standard Python environment. The planned dashboard shall run in current web browsers.
   \item \textbf{Interoperability:} Module outputs shall use documented, structured formats such as JSON or CSV, so that they can be consumed by the planned backend or exported to other tools.
   \item \textbf{Availability:} In the demonstration environment, the integrated system shall remain usable if one module fails, reporting that module as unavailable instead of stopping completely.
   \item \textbf{Reproducibility:} Data splits, random seeds and configurations shall be fixed and saved, so that every reported result can be regenerated.
\end{itemize}

\subsection{Performance Requirements}

\begin{itemize}
   \setlength{\itemsep}{1pt}
   \setlength{\parskip}{0pt}
   \setlength{\parsep}{0pt}
   \item \textbf{Detection quality.} On its benchmark test data, each supervised classifier shall reach a macro F1-score of at least 0.85. The current values are 0.958 for network threat classification and 0.877 for the malware category stage. The anomaly detector shall reach an F1-score of at least 0.80 at its calibrated threshold; the current mean values are 0.835 on CIC-IDS2017 and 0.899 on NSL-KDD.
   \item \textbf{False alarms.} The anomaly detector's false-positive rate shall be configurable through its threshold, and shall be reported on every evaluation. A target of at most 5\% on traffic from the deployment network shall be used when recalibrating. The current test values are 2.9\% on CIC-IDS2017 and 11.9\% on NSL-KDD, so the latter needs recalibration.
   \item \textbf{Batch throughput.} Scoring a batch of 1,000 flow records with a trained model shall take less than one second on the development machine. The anomaly module currently takes about 35--52 milliseconds per 1,000 flows on a 12-core machine. Timings for the other modules will be measured during integration.
   \item \textbf{Training time.} Retraining a module on its benchmark dataset shall be possible on a standard development machine without specialised hardware.
   \item \textbf{Responsiveness.} The planned dashboard shall display an analysed batch of up to 10,000 records within a few seconds of the analysis completing.
\end{itemize}

\subsection{Safety Requirements}

\begin{itemize}
   \setlength{\itemsep}{1pt}
   \setlength{\parskip}{0pt}
   \setlength{\parsep}{0pt}
   \item The system shall not execute blocking, isolation or other countermeasures automatically. Recommended responses shall be presented to an analyst for decision, because a wrong automatic action could disrupt legitimate services.
   \item The system shall be developed and tested only in a controlled environment, using public datasets and systems for which permission exists. It shall not be used to scan or attack third-party systems.
   \item The project shall not create, modify or execute malware. The malware module works only on numerical feature vectors from a public research dataset.
   \item When a module receives input it cannot process correctly (for example, missing feature columns), it shall report an error or warning instead of returning results that look normal (see FR-M2 and FR-A8).
   \item The limitations of each model, such as benchmark-only evaluation and weak results on rare classes, shall be shown alongside its results so that its outputs are not over-trusted.
\end{itemize}

\subsection{Other Non-Functional Requirements}

\begin{itemize}
   \setlength{\itemsep}{1pt}
   \setlength{\parskip}{0pt}
   \setlength{\parsep}{0pt}
   \item \textbf{Security of the platform:} The planned backend and dashboard shall require authentication, restrict stored security data by role, and protect it in storage and in transit.
   \item \textbf{Privacy:} Direct identifiers such as IP addresses shall not be used as model inputs. They may appear only in alerts, as context for the analyst.
   \item \textbf{Legal and ethical use:} Datasets shall be used according to their published terms, and all sources shall be cited.
   \item \textbf{Portability:} The trained models and their preprocessing objects shall be stored together, so that a module can be moved to another machine and loaded without retraining.
\end{itemize}
